\chapter{Values, Types, and the Shape of Things}\label{ch:values-types}

\index{types}\index{values}

Before you can freeze something, you need something \emph{to} freeze.
Before you can match a pattern, you need data shaped in a way the pattern can describe.
Everything in Lattice starts with values---the fundamental atoms of computation---and their types.

In this chapter, we will meet every core value type in Lattice, learn the operators that transform them, and discover the surprisingly rich world of string literals.
By the end, you will have a complete mental model of the raw materials Lattice gives you to work with.

%% ───────────────────────────────────────────────────────────
\section{Integers}\label{sec:integers}
\index{Int}\index{integers}

Integers in Lattice are 64-bit signed values, stored internally as C's \lstinline|int64_t| (you can verify this in \filename{include/value.h}, where \lstinline|VAL_INT| maps to the \lstinline|int_val| field of the value union).
This gives you a range from $-2^{63}$ to $2^{63}-1$---large enough for most practical purposes.

\begin{lstlisting}[language=Lattice, caption={Integer basics}]
let population = 8_100_000_000
let negative = -42
let zero = 0

print(typeof(population))  // Int
print(typeof(zero))        // Int
\end{lstlisting}

\begin{defnbox}{Integer Range}
Lattice integers are 64-bit signed: they range from $-9{,}223{,}372{,}036{,}854{,}775{,}808$ to $9{,}223{,}372{,}036{,}854{,}775{,}807$.
There is no automatic promotion to big integers---if you overflow, you'll get unexpected results.
For most applications, 64 bits is more than sufficient.
\end{defnbox}

%% ───────────────────────────────────────────────────────────
\section{Floats}\label{sec:floats}
\index{Float}\index{floats}

Floating-point numbers use 64-bit double precision (C's \lstinline|double|).
Any number literal containing a decimal point is a float:

\begin{lstlisting}[language=Lattice, caption={Float basics}]
let pi = 3.14159
let temperature = -40.0
let ratio = 0.618

print(typeof(pi))           // Float
print(typeof(temperature))  // Float
\end{lstlisting}

\begin{warnbox}{Integer Division vs.\ Float Division}
Division between two integers produces an integer (truncated toward zero).
If you want a decimal result, make sure at least one operand is a float:
\begin{lstlisting}[language=Lattice, numbers=none]
print(7 / 2)       // 3 (integer division)
print(7.0 / 2.0)   // 3.5 (float division)
print(7 / 2.0)     // 3.5 (mixed: promotes to float)
\end{lstlisting}
\end{warnbox}

Lattice provides two functions for inspecting special float values:

\begin{lstlisting}[language=Lattice, caption={Special float values}]
print(is_nan(0.0 / 0.0))   // true
print(is_inf(1.0 / 0.0))   // true
\end{lstlisting}

You can convert between integers and floats explicitly:

\begin{lstlisting}[language=Lattice, caption={Numeric conversions}]
let x = to_float(42)     // 42.0
let y = to_int(3.14)     // 3 (truncates toward zero)
let z = round(3.7)       // 4
let w = floor(3.9)       // 3
let v = ceil(3.1)        // 4
\end{lstlisting}

%% ───────────────────────────────────────────────────────────
\section{Booleans}\label{sec:booleans}
\index{Bool}\index{booleans}

Booleans are \lstinline|true|\index{true} or \lstinline|false|\index{false}.
No surprises here:

\begin{lstlisting}[language=Lattice, caption={Boolean basics}]
let is_valid = true
let has_errors = false

print(typeof(is_valid))  // Bool
print(!is_valid)         // false
print(is_valid && has_errors)  // false
print(is_valid || has_errors)  // true
\end{lstlisting}

\subsection{Truthiness}\label{subsec:truthiness}
\index{truthiness}

In contexts where a boolean is expected (like \lstinline|if| conditions), Lattice applies \emph{truthiness} rules:

\begin{itemize}
  \item \lstinline|false| is falsy.
  \item \lstinline|nil| is falsy.
  \item \lstinline|0| (integer zero) is falsy.
  \item \lstinline|""| (empty string) is falsy.
  \item Everything else is truthy---including empty arrays and empty maps.
\end{itemize}

\begin{lstlisting}[language=Lattice, caption={Truthiness in action}]
if 42 { print("truthy") }         // truthy
if "" { print("truthy") } else { print("falsy") }  // falsy
if nil { print("truthy") } else { print("falsy") }  // falsy
if [] { print("truthy") }         // truthy (empty array is truthy!)
\end{lstlisting}

\begin{warnbox}{Empty Arrays Are Truthy}
Unlike Python, empty arrays and empty maps are \emph{truthy} in Lattice.
If you want to check whether a collection is empty, use \lstinline|.len() == 0| or \lstinline|.is_empty()| (for strings).
\end{warnbox}

%% ───────────────────────────────────────────────────────────
\section{Strings}\label{sec:strings}
\index{String}\index{strings}

Strings in Lattice are immutable, UTF-8 encoded sequences of characters.
They are the most feature-rich of the primitive types, with three different literal syntaxes and over 20 built-in methods.

\subsection{Double-Quoted Strings}\label{subsec:double-strings}
\index{strings!double-quoted}

The standard string literal uses double quotes and supports both escape sequences and interpolation:

\begin{lstlisting}[language=Lattice, caption={Double-quoted strings}]
let greeting = "Hello, World!"
let with_newline = "line one\nline two"
let with_tab = "column A\tcolumn B"
\end{lstlisting}

The full set of escape sequences is:

\begin{table}[h]
\centering
\begin{tabular}{ll}
\toprule
\textbf{Escape} & \textbf{Character} \\
\midrule
\lstinline|\\n| & Newline \\
\lstinline|\\t| & Tab \\
\lstinline|\\r| & Carriage return \\
\lstinline|\\0| & Null byte \\
\lstinline|\\\\| & Literal backslash \\
\lstinline|\\"| & Literal double quote \\
\lstinline|\\'| & Literal single quote \\
\lstinline|\\$| & Literal dollar sign (prevents interpolation) \\
\lstinline|\\xHH| & Hex byte (e.g., \lstinline|\\x41| for 'A') \\
\bottomrule
\end{tabular}
\caption{String escape sequences}\label{tab:escape-sequences}
\index{escape sequences}
\end{table}

\subsection{String Interpolation}\label{subsec:string-interpolation}
\index{string interpolation}

Inside double-quoted strings, \lstinline|${expr}| embeds the result of any expression:

\begin{lstlisting}[language=Lattice, caption={String interpolation}]
let language = "Lattice"
let major = 0
let minor = 3

print("Welcome to ${language}")
// Welcome to Lattice

print("Version ${major}.${minor}")
// Version 0.3

print("2 + 2 = ${2 + 2}")
// 2 + 2 = 4

let name = "world"
print("${"hello".to_upper()}, ${name}!")
// HELLO, world!
\end{lstlisting}

The expression inside the braces can be arbitrarily complex---function calls, method chains, arithmetic, even nested strings.
The lexer handles this by tracking brace depth, recursively lexing the expression tokens embedded within the string (you can see this machinery in \filename{src/lexer.c}, where interpolated strings are split into \lstinline|TOK_INTERP_START|, expression tokens, and \lstinline|TOK_INTERP_END| segments).

To include a literal \lstinline|${| in a string, escape the dollar sign:

\begin{lstlisting}[language=Lattice, caption={Escaping interpolation}]
print("Use \${expr} for interpolation")
// Use ${expr} for interpolation
\end{lstlisting}

\subsection{Single-Quoted Strings}\label{subsec:single-strings}
\index{strings!single-quoted}

Single-quoted strings have \textbf{no interpolation}\index{string interpolation!disabling}.
The \lstinline|${...}| syntax is treated as literal characters:

\begin{lstlisting}[language=Lattice, caption={Single-quoted strings: no interpolation}]
let pattern = 'Hello, ${name}'
print(pattern)
// Hello, ${name}
\end{lstlisting}

Single-quoted strings still support escape sequences (\lstinline|\\n|, \lstinline|\\t|, etc.), so they are not fully ``raw'' strings.
Their primary use case is when you have strings that contain dollar signs and braces and you don't want to worry about accidental interpolation:

\begin{lstlisting}[language=Lattice, caption={Single-quoted strings for templates}]
let template = 'Dear ${recipient}, your order #${order_id} is ready.'
print(template)
// Dear ${recipient}, your order #${order_id} is ready.
\end{lstlisting}

\begin{tipbox}{When to Use Single vs.\ Double Quotes}
Use double quotes when you want interpolation: \lstinline|"Hello, ${name}"|.
Use single quotes when you're writing templates, regex patterns, or any string where \lstinline|$\{...\}| should be literal text.
\end{tipbox}

\subsection{Triple-Quoted Strings}\label{subsec:triple-strings}
\index{strings!triple-quoted}\index{multi-line strings}

Triple-quoted strings (\lstinline|""" ... """|) span multiple lines and support automatic dedenting:

\begin{lstlisting}[language=Lattice, caption={Triple-quoted strings}]
let html = """
    <div>
        <h1>Welcome</h1>
        <p>Hello, world!</p>
    </div>
    """
print(html)
\end{lstlisting}

Output (note how the leading whitespace has been stripped):

\begin{lstlisting}[numbers=none]
<div>
    <h1>Welcome</h1>
    <p>Hello, world!</p>
</div>
\end{lstlisting}

The dedenting algorithm uses the indentation of the closing \lstinline|"""| as the baseline and strips that much leading whitespace from every line.
This lets you indent your multi-line strings to match the surrounding code without polluting the string content.

Triple-quoted strings support interpolation:

\begin{lstlisting}[language=Lattice, caption={Triple-quoted strings with interpolation}]
let name = "Lattice"
let version = "0.3.28"
let banner = """
    ============================
     ${name} v${version}
     A crystallization language
    ============================
    """
print(banner)
\end{lstlisting}

\subsection{Common String Methods}\label{subsec:string-methods}
\index{strings!methods}

Strings come with a rich set of methods.
Here are some of the most commonly used ones:

\begin{lstlisting}[language=Lattice, caption={String methods}]
let message = "  Hello, Lattice!  "

print(message.trim())            // "Hello, Lattice!"
print(message.len())             // 20
print(message.contains("Lat"))   // true
print(message.to_upper())        // "  HELLO, LATTICE!  "
print(message.replace("Lattice", "World"))  // "  Hello, World!  "
print(message.split(","))        // ["  Hello", " Lattice!  "]

let path = "/home/user/documents/report.txt"
print(path.starts_with("/home"))  // true
print(path.ends_with(".txt"))     // true
print(path.index_of("user"))     // 6
\end{lstlisting}

For a complete reference of string methods, see Chapter~8 (\emph{Strings in Depth}).

%% ───────────────────────────────────────────────────────────
\section{Nil and Unit}\label{sec:nil-unit}

Lattice has two ``empty'' types, and understanding the difference matters.

\subsection{Nil}\label{subsec:nil}
\index{nil}\index{Nil}

\lstinline|nil| represents ``no value here''---the explicit absence of a meaningful value.
It's the value returned by \lstinline|Map.get()| for a missing key, and it's what you use when you want to intentionally store ``nothing'':

\begin{lstlisting}[language=Lattice, caption={Nil basics}]
let nothing = nil
print(typeof(nothing))  // Nil

flux m = Map::new()
m["key"] = "value"
print(m.get("key"))      // value
print(m.get("missing"))  // nil
\end{lstlisting}

\lstinline|nil| is \emph{falsy}---it evaluates to false in boolean contexts.
This makes it natural to use with the nil coalescing operator\index{nil coalescing}:

\begin{lstlisting}[language=Lattice, caption={Nil coalescing with ??}]
let port = nil ?? 8080
print(port)  // 8080

let name = "Alice" ?? "Anonymous"
print(name)  // Alice (left side is not nil, so it wins)
\end{lstlisting}

\subsection{Unit}\label{subsec:unit}
\index{unit}\index{Unit}

\lstinline|unit| is the type of ``nothing happened''---the return type of statements, \lstinline|print()| calls, and variable assignments.
You rarely encounter it explicitly, but it's always there behind the scenes:

\begin{lstlisting}[language=Lattice, caption={Unit in action}]
let result = print("hello")
print(typeof(result))  // Unit
\end{lstlisting}

In the REPL, both \lstinline|nil| and \lstinline|unit| are suppressed---neither produces the \lstinline|=>| output line.

\begin{defnbox}{Nil vs.\ Unit}
\lstinline|nil| means ``this value is intentionally empty.''
It can be stored in variables, returned from functions, and compared with \lstinline|==|.

\lstinline|unit| means ``this expression produces no value.''
It's the implicit return of statements and void-like functions.
You typically don't store it or test for it.

Think of \lstinline|nil| as an empty box and \lstinline|unit| as the absence of a box.
\end{defnbox}

%% ───────────────────────────────────────────────────────────
\section{Type Checking with \texttt{typeof()}}\label{sec:typeof}
\index{typeof()}

The \lstinline|typeof()| built-in function returns a string naming the type of any value:

\begin{lstlisting}[language=Lattice, caption={typeof() for every core type}]
print(typeof(42))             // Int
print(typeof(3.14))           // Float
print(typeof(true))           // Bool
print(typeof("hello"))        // String
print(typeof(nil))            // Nil
print(typeof([1, 2, 3]))     // Array
print(typeof(Map::new()))    // Map
print(typeof(Set::new()))    // Set
print(typeof((1, 2)))        // Tuple
print(typeof(0..10))         // Range
print(typeof(|x| x))         // Closure
\end{lstlisting}

\lstinline|typeof()| returns the string name, not a type object.
You compare it with string equality:

\begin{lstlisting}[language=Lattice, caption={Runtime type checking}]
fn describe(value: any) -> String {
    if typeof(value) == "Int" {
        "an integer: ${to_string(value)}"
    } else if typeof(value) == "String" {
        "a string: ${value}"
    } else if typeof(value) == "Array" {
        "an array of length ${to_string(value.len())}"
    } else {
        "something else: ${to_string(value)}"
    }
}

print(describe(42))          // an integer: 42
print(describe("hello"))     // a string: hello
print(describe([1, 2, 3]))  // an array of length 3
\end{lstlisting}

For struct instances, \lstinline|typeof()| returns the struct's name:

\begin{lstlisting}[language=Lattice, caption={typeof() with structs}]
struct User { name: String, age: Int }
let alice = User { name: "Alice", age: 30 }
print(typeof(alice))  // User
\end{lstlisting}

\begin{tipbox}{Type Annotations vs.\ typeof()}
Function parameter type annotations (\lstinline|fn foo(x: Int)|) and \lstinline|typeof()| serve different purposes.
Annotations are checked automatically on every call and produce clear error messages.
\lstinline|typeof()| is for when you need to branch on type at runtime---rare in well-typed code, but useful for utility functions, serializers, and debugging.
\end{tipbox}

%% ───────────────────────────────────────────────────────────
\section{Number Literals in Detail}\label{sec:number-literals}
\index{number literals}

Lattice supports several formats for writing numeric literals, all designed for readability.

\subsection{Underscore Separators}\label{subsec:underscores}
\index{number literals!underscores}

You can place underscores between digits in both integer and float literals.
The underscores are purely visual---they are stripped before the number is parsed:

\begin{lstlisting}[language=Lattice, caption={Underscore separators}]
let population = 8_100_000_000
let bytes = 1_048_576
let precise_pi = 3.14_159_265

print(population)   // 8100000000
print(bytes)        // 1048576
print(precise_pi)   // 3.14159265
\end{lstlisting}

This is especially useful for large numbers where counting zeros is error-prone.

\subsection{Hexadecimal Literals}\label{subsec:hex-literals}
\index{number literals!hexadecimal}

Integer literals can be written in hexadecimal with the \lstinline|0x| prefix:

\begin{lstlisting}[language=Lattice, caption={Hexadecimal literals}]
let red = 0xFF0000
let mask = 0xDEAD_BEEF
let byte = 0x7F

print(red)    // 16711680
print(mask)   // 3735928559
print(byte)   // 127
\end{lstlisting}

Hex digits can be uppercase or lowercase, and underscore separators work with hex too.
The lexer (in \filename{src/lexer.c}) detects the \lstinline|0x| or \lstinline|0X| prefix, reads hex digits with \lstinline|isxdigit()|, strips underscores, and converts using \lstinline|strtoll()| with base 16.

\subsection{No Scientific Notation}\label{subsec:no-scientific}

Lattice does not currently support scientific notation (like \lstinline|1.5e10| or \lstinline|3e-4|).
If you need to express very large or very small numbers, use \lstinline|pow()|:

\begin{lstlisting}[language=Lattice, caption={Expressing large numbers}]
let avogadro = 6.022 * pow(10.0, 23.0)
let planck = 6.626 * pow(10.0, -34.0)
\end{lstlisting}

%% ───────────────────────────────────────────────────────────
\section{Operators}\label{sec:operators}
\index{operators}

Lattice provides a comprehensive set of operators that will feel familiar to anyone coming from C, JavaScript, or Python.

\subsection{Arithmetic Operators}\label{subsec:arithmetic}
\index{operators!arithmetic}

\begin{table}[h]
\centering
\begin{tabular}{llll}
\toprule
\textbf{Operator} & \textbf{Description} & \textbf{Example} & \textbf{Result} \\
\midrule
\lstinline|+| & Addition & \lstinline|3 + 4| & \lstinline|7| \\
\lstinline|-| & Subtraction & \lstinline|10 - 3| & \lstinline|7| \\
\lstinline|*| & Multiplication & \lstinline|6 * 7| & \lstinline|42| \\
\lstinline|/| & Division & \lstinline|15 / 4| & \lstinline|3| \\
\lstinline|%| & Modulo & \lstinline|17 % 5| & \lstinline|2| \\
\bottomrule
\end{tabular}
\caption{Arithmetic operators}\label{tab:arithmetic-operators}
\end{table}

The \lstinline|+| operator also works for string concatenation\index{string concatenation}:

\begin{lstlisting}[language=Lattice, caption={String concatenation with +}]
let full = "Hello" + ", " + "World!"
print(full)  // Hello, World!
\end{lstlisting}

\subsection{Comparison Operators}\label{subsec:comparison}
\index{operators!comparison}

\begin{table}[h]
\centering
\begin{tabular}{llll}
\toprule
\textbf{Operator} & \textbf{Description} & \textbf{Example} & \textbf{Result} \\
\midrule
\lstinline|==| & Equal & \lstinline|3 == 3| & \lstinline|true| \\
\lstinline|!=| & Not equal & \lstinline|3 != 4| & \lstinline|true| \\
\lstinline|<|  & Less than & \lstinline|3 < 5| & \lstinline|true| \\
\lstinline|>|  & Greater than & \lstinline|5 > 3| & \lstinline|true| \\
\lstinline|<=| & Less or equal & \lstinline|3 <= 3| & \lstinline|true| \\
\lstinline|>=| & Greater or equal & \lstinline|4 >= 5| & \lstinline|false| \\
\bottomrule
\end{tabular}
\caption{Comparison operators}\label{tab:comparison-operators}
\end{table}

Comparison operators work on integers, floats, strings (lexicographic order), and booleans.
Equality (\lstinline|==|) uses structural comparison for arrays, maps, and structs:

\begin{lstlisting}[language=Lattice, caption={Structural equality}]
print([1, 2, 3] == [1, 2, 3])   // true
print([1, 2] == [1, 2, 3])      // false
print("abc" < "abd")             // true (lexicographic)
\end{lstlisting}

\subsection{Logical Operators}\label{subsec:logical}
\index{operators!logical}

\begin{table}[h]
\centering
\begin{tabular}{llll}
\toprule
\textbf{Operator} & \textbf{Description} & \textbf{Example} & \textbf{Result} \\
\midrule
\lstinline|&&| & Logical AND & \lstinline|true && false| & \lstinline|false| \\
\lstinline!||! & Logical OR  & \lstinline!true || false! & \lstinline|true| \\
\lstinline|!|  & Logical NOT & \lstinline|!true| & \lstinline|false| \\
\bottomrule
\end{tabular}
\caption{Logical operators}\label{tab:logical-operators}
\end{table}

Both \lstinline|&&| and \lstinline!||! are short-circuiting\index{short-circuit evaluation}: the right operand is only evaluated if needed.

\begin{lstlisting}[language=Lattice, caption={Short-circuit evaluation}]
// The second condition is never evaluated
let safe = false && (1 / 0 > 0)
print(safe)  // false

// The second condition is never evaluated
let ok = true || (1 / 0 > 0)
print(ok)    // true
\end{lstlisting}

\subsection{Bitwise Operators}\label{subsec:bitwise}
\index{operators!bitwise}

Lattice provides the full set of bitwise operations on integers:

\begin{table}[h]
\centering
\begin{tabular}{llll}
\toprule
\textbf{Operator} & \textbf{Description} & \textbf{Example} & \textbf{Result} \\
\midrule
\lstinline|&| & Bitwise AND & \lstinline|0xFF & 0x0F| & \lstinline|15| \\
\lstinline!|! & Bitwise OR  & \lstinline!0xF0 | 0x0F! & \lstinline|255| \\
\lstinline|^| & Bitwise XOR & \lstinline|0xFF ^ 0x0F| & \lstinline|240| \\
\lstinline|~| & Bitwise NOT & \lstinline|~0| & \lstinline|-1| \\
\lstinline|<<| & Left shift  & \lstinline|1 << 8| & \lstinline|256| \\
\lstinline|>>| & Right shift & \lstinline|256 >> 4| & \lstinline|16| \\
\bottomrule
\end{tabular}
\caption{Bitwise operators}\label{tab:bitwise-operators}
\end{table}

Bitwise operators are essential for systems programming tasks like flag manipulation, color encoding, and protocol parsing:

\begin{lstlisting}[language=Lattice, caption={Bitwise operations in practice}]
// Extract RGB components from a 24-bit color
let color = 0xFF8040
let red   = (color >> 16) & 0xFF   // 255
let green = (color >> 8) & 0xFF    // 128
let blue  = color & 0xFF           // 64

print("R=${to_string(red)} G=${to_string(green)} B=${to_string(blue)}")
// R=255 G=128 B=64
\end{lstlisting}

\subsection{Compound Assignment}\label{subsec:compound-assignment}
\index{operators!compound assignment}

Every arithmetic and bitwise operator has a compound assignment form:

\begin{lstlisting}[language=Lattice, caption={Compound assignment operators}]
flux score = 100
score += 10    // score = 110
score -= 5     // score = 105
score *= 2     // score = 210
score /= 3     // score = 70
score %= 30    // score = 10

flux flags = 0xFF
flags &= 0x0F   // flags = 15
flags |= 0xF0   // flags = 255
flags ^= 0x0F   // flags = 240
flags <<= 4     // flags = 3840
flags >>= 8     // flags = 15
\end{lstlisting}

\begin{notebox}{Compound Assignment Requires \texttt{flux}}
Compound assignment operators modify the variable in place.
This only works with \lstinline|flux| (mutable) variables.
Attempting to use \lstinline|+=| on a \lstinline|let| or \lstinline|fix| binding that refers to a crystal value will produce an error.
\end{notebox}

\subsection{The Nil Coalescing Operator}\label{subsec:nil-coalescing}
\index{operators!nil coalescing}\index{??}

The \lstinline|??| operator returns the left side if it is not \lstinline|nil|; otherwise, it returns the right side:

\begin{lstlisting}[language=Lattice, caption={Nil coalescing}]
let name = nil ?? "Anonymous"
print(name)  // Anonymous

let title = "Engineer" ?? "Unknown"
print(title)  // Engineer

// Chains: first non-nil wins
let result = nil ?? nil ?? "found it"
print(result)  // found it
\end{lstlisting}

This is particularly useful when reading from maps, where missing keys return \lstinline|nil|:

\begin{lstlisting}[language=Lattice, caption={Nil coalescing with map lookups}]
flux config = Map::new()
config["host"] = "localhost"

let host = config.get("host") ?? "0.0.0.0"
let port = config.get("port") ?? 8080

print(host)  // localhost
print(port)  // 8080
\end{lstlisting}

\subsection{Optional Chaining}\label{subsec:optional-chaining}
\index{operators!optional chaining}\index{?.}

The \lstinline|?.|\index{?.} operator safely navigates through potentially \lstinline|nil| values.
If the receiver is \lstinline|nil|, the entire chain short-circuits to \lstinline|nil| instead of producing an error:

\begin{lstlisting}[language=Lattice, caption={Optional chaining}]
let user = nil
print(user?.name)           // nil (no error)
print(user?.address?.city)  // nil (chain short-circuits)

// Combine with ?? for defaults
let city = user?.address?.city ?? "Unknown"
print(city)  // Unknown
\end{lstlisting}

\subsection{The Range Operator}\label{subsec:range-operator}
\index{operators!range}\index{..}

The \lstinline|..| operator creates a range---a pair of integers representing a half-open interval (\lstinline|start| inclusive, \lstinline|end| exclusive):

\begin{lstlisting}[language=Lattice, caption={Range operator}]
for i in 0..5 {
    print(i)  // 0, 1, 2, 3, 4
}

print(typeof(0..10))  // Range
\end{lstlisting}

Ranges are used with \lstinline|for| loops and in pattern matching.

%% ───────────────────────────────────────────────────────────
\section{Comments}\label{sec:comments}
\index{comments}

Lattice supports three styles of comments.

\subsection{Line Comments}\label{subsec:line-comments}
\index{comments!line}

A \lstinline|//| starts a comment that extends to the end of the line:

\begin{lstlisting}[language=Lattice, caption={Line comments}]
let x = 42  // this is a comment
// this entire line is a comment
\end{lstlisting}

\subsection{Block Comments}\label{subsec:block-comments}
\index{comments!block}

Block comments use \lstinline|/* ... */| and can span multiple lines.
They can be \textbf{nested}\index{comments!nested}, which is particularly useful for commenting out code that already contains comments:

\begin{lstlisting}[language=Lattice, caption={Block comments (nestable)}]
/* This is a block comment */

/*
 * Multi-line block comment.
 * Often used for function documentation.
 */

/* You can /* nest */ block comments */
\end{lstlisting}

The nesting support is implemented by tracking a depth counter in the lexer---each \lstinline|/*| increments it and each \lstinline|*/| decrements it.
The comment only ends when the depth reaches zero.

\subsection{Doc Comments}\label{subsec:doc-comments}
\index{comments!doc}\index{doc comments}

Lines starting with \lstinline|///| are \emph{doc comments}.
They are treated as regular comments by the interpreter, but the \lstinline|clat doc| tool extracts them to generate documentation:

\begin{lstlisting}[language=Lattice, caption={Doc comments}]
/// Calculate the area of a circle.
/// Takes a radius and returns the area.
fn circle_area(radius: Float) -> Float {
    3.14159 * radius * radius
}
\end{lstlisting}

We'll cover the documentation generator in detail in Chapter~30 (\emph{The Formatter and Doc Generator}).
For now, get in the habit of documenting your functions with \lstinline|///|---your future self will thank you.

%% ───────────────────────────────────────────────────────────
\section{Putting It All Together}\label{sec:types-together}

Let's write a small program that uses several of the types and operators we've covered:

\begin{lstlisting}[language=Lattice, caption={A temperature converter using multiple types}]
/// Convert a temperature between Fahrenheit and Celsius.
/// The `from` parameter must be "F" or "C".
fn convert_temp(value: Float, from: String) -> String {
    let result = if from == "F" {
        (value - 32.0) * 5.0 / 9.0
    } else {
        value * 9.0 / 5.0 + 32.0
    }

    let to_unit = if from == "F" { "C" } else { "F" }
    let rounded = round(result * 100.0) / 100.0

    "${to_string(value)} ${from} = ${to_string(rounded)} ${to_unit}"
}

// Build a conversion table
let temps = [0.0, 32.0, 72.0, 100.0, 212.0]
for temp in temps {
    print(convert_temp(temp, "F"))
}
/*
    0.0 F = -17.78 C
    32.0 F = 0.0 C
    72.0 F = 22.22 C
    100.0 F = 37.78 C
    212.0 F = 100.0 C
*/
\end{lstlisting}

This program demonstrates:
\begin{itemize}
  \item \textbf{Floats} for temperature values
  \item \textbf{Strings} with interpolation for formatted output
  \item \textbf{Comparison operators} (\lstinline|==|) for branching
  \item \textbf{Arithmetic operators} for the conversion formula
  \item \textbf{Arrays} to hold a list of temperatures
  \item \textbf{Doc comments} on the function
  \item \textbf{Block comments} for expected output
  \item The \lstinline|if|/\lstinline|else| expression returning a value
\end{itemize}

%% ───────────────────────────────────────────────────────────
\section{Exercises}\label{sec:ch03-exercises}

\begin{enumerate}
  \item \textbf{Type Explorer.}
  In the REPL, use \lstinline|typeof()| on at least 10 different values, including an array of arrays, a map with string keys, a range, a closure, and \lstinline|nil|.
  Write down any surprises.

  \item \textbf{Hex Color Parser.}
  Write a function \lstinline|fn parse_color(hex: Int) -> [Int]| that takes a hex color like \lstinline|0xFF8040| and returns an array \lstinline|[r, g, b]| using bitwise operators.
  Test it with several colors.

  \item \textbf{String Processing.}
  Write a function that takes a full name like \lstinline|"ada lovelace"| and returns it in title case: \lstinline|"Ada Lovelace"|.
  Use string methods.
  (Hint: \lstinline|.title_case()| exists!)

  \item \textbf{Nil Safety.}
  Write a function \lstinline|fn safe_get(m: Map, key: String, default_val: any) -> any| that returns the map's value for the given key, or \lstinline|default_val| if the key is missing.
  Use the \lstinline|??| operator.

  \item \textbf{Triple-Quote Art.}
  Use a triple-quoted string with interpolation to create an ASCII art box around a message.
  The function should take a message string and return the boxed version with proper padding.
\end{enumerate}

%% ───────────────────────────────────────────────────────────
\section*{What's Next}\label{sec:ch03-whats-next}

We've surveyed the raw materials---integers, floats, booleans, strings, nil, and unit---and the tools for working with them.
But so far, our values have been rootless: created, used, and forgotten.
In the next chapter, we will learn how to \emph{bind} values to names with \lstinline|let|, \lstinline|flux|, and \lstinline|fix|, and we'll get our first real look at the phase system that gives Lattice its crystalline character.
